% !TeX root = ../paper.tex

\section{引言}

蛋白质功能由结构与构象运动共同决定。实验结构和结构预测模型通常给出单个或少量
构象，而真实蛋白质在热涨落下持续采样构象集合；界面形成还会重新分配局部运动，
使部分残基受限，并可能在界面外产生协同或补偿性变化。分子动力学（molecular
dynamics，MD）能够在原子分辨率下刻画这些过程，但其计算成本限制了大规模应用。
ATLAS 和 DynaRepo 等标准化轨迹资源缓解了动态数据分散、模拟协议不一致和重复不足
的问题\cite{vandermeersche2024atlas,mokhtari2026dynarepo}，同时也使“由一个静态
结构预测残基级动态量”成为可监督学习任务。本文关注三个互补描述符：均方根波动
（root mean square fluctuation，RMSF）度量主链位置涨落，侧链构象熵
$H_\chi$ 描述旋转异构状态的多样性，溶剂可及表面积（solvent-accessible surface
area，SASA）反映局部溶剂暴露。三者均可由轨迹估计，但分别对应不同层面的动态
信息，不能相互替代。

从静态结构预测柔性的既有方法大体分为物理近似、结构统计和数据驱动模型。正则模态
与粗粒化模拟以较低成本近似集体运动；FastProtFlex 使用原子邻接图的局部拓扑统计
预测相对 RMSF\cite{praznikar2026fast}，ProDy 提供基于 C$\alpha$ 接触网络的
ANM/GNM 模态分析\cite{zhang2021prody}；RMSF-net 将 PDB 或 cryo-EM 密度体素化后
用三维卷积网络预测柔性\cite{song2024rmsfnet}；Flexpert 与 BackFlip 分别以
序列--结构表示和等变主链网络学习残基级柔性
\cite{kouba2025flexpert,viliuga2025backflip}。与这些直接预测器不同，AlphaFlow
通过 flow matching 生成构象集合，可从采样坐标导出 RMSF\cite{jing2024alphaflow}。
另一方面，几何深度学习通过距离、局部坐标系或等变表示消除任意全局坐标系的影响。
SE(3)-Transformer 将旋转和平移对称性纳入注意力计算\cite{fuchs2020se3}，
AlphaFold2 的不变点注意力（invariant point attention，IPA）利用局部框架和三维
查询点更新残基表示\cite{jumper2021alphafold}；
PeSTo 与 DeepInteract 进一步表明，几何注意力能够直接建模蛋白质界面及伙伴特异的
残基关系\cite{krapp2023pesto,morehead2022deepinteract}。这些进展说明静态结构中
包含可学习的柔性和界面线索，但并不自动保证模型能够恢复结合伙伴所对应的动力学
变化。

现有评估通常把一个蛋白质内部的预测 RMSF 与参考 RMSF 相关联。该指标适合衡量
“哪些残基相对更柔”，却不能回答两个与蛋白质结合直接相关的问题。第一，模型在
完整复合物上取得较高的总体相关性，可能主要来自稳定的二级结构和非界面区域，
而非对界面动力学的正确刻画。第二，将伙伴删除所造成的预测变化同时包含跨链接触、
输入长度和注意力归一化变化；如果两种输入仍共享结合态坐标，该差值只能称为模型
上下文响应，不能等同于具有独立游离态轨迹支持的 apo/holo 响应。抗体研究中对
结合前后 CDR 构象变化的分析同样强调，局部构象改变、相对框架位移和构象选择必须
分别定义\cite{liu2024cdr}。因此，柔性模型需要同时接受整体排序、界面分层、残基级
定位和受控上下文消融，而不能只报告一个汇总相关系数。

为此，本文提出 FlexPPI。模型以氨基酸类别、实验条件和主链坐标为输入，用固定
N--C$\alpha$--C 局部框架构造旋转和平移不变的残基对几何，并通过四层 IPA 在完整
多链结构中传播伙伴信息。概率状态头为每个残基预测 RMSF、$H_\chi$ 和 SASA 的均值
与方差；噪声感知目标将重复轨迹或时间窗口间方差显式并入训练。模型首先在 ATLAS
和 DynaRepo 上学习绝对状态。训练结束后，我们冻结检查点，在 ITrace 的 245 个
TCR--pMHC 复合体和 735 条轨迹上构建外部评估。主要伙伴消融保持结合态坐标不变，
仅删除对侧组分：pMHC 残基比较完整五链与 pMHC-only 输入，TCR 残基比较完整五链与
TCR$\alpha\beta$-only 输入；五条孤立链只作为进一步拆解子复合物的敏感性分析。
该设计将“静态结构能否预测 MD 真值”与“模型是否响应对侧伙伴上下文”分开检验，
并避免把 pMHC 内部或 TCR$\alpha\beta$ 内部相互作用混入主要伙伴效应。

外部结果揭示了模型能力的明确边界。排除唯一精确 PDB 重叠 3KXF 后，FlexPPI 在
244 条抗原肽内的平均残基 Spearman 为 $0.302$，能够恢复九肽 P1 活跃和 P2/P8
受限的群体趋势，但遗漏 P6 活动峰；最活跃残基命中率仅为 34.0\%。在全部链上，
完整复合物预测的平均 Spearman 为 $0.740$，高于逐链 BackFlip-seq 的 $0.704$、
FastProtFlex 的 $0.685$ 和 RMSF-net-PDB 的 $0.526$；但是，BackFlip-seq 在线性
相关和标准化误差上更好，结构版 BackFlip 在 CDR3 内也明显领先。FlexPPI 的总体
排序增益主要来自非界面区域。进一步的固定坐标伙伴消融显示，完整输入相对删伙伴
子复合物将 Spearman 从 0.691 提高到 0.740，并将绝对 MAE 从
$0.373\,\text{\AA}$ 降至 $0.346\,\text{\AA}$；但加入伙伴会近乎单向地下调接触
残基的 RMSF、$H_\chi$ 和 SASA。该偏置改善了 MHC 接触面的相对刚性，却将 MD 中
较柔的 TCR 接触面和 CDR1/2$\beta$ 错误压低。模型由此学到了可定位但过度单向的
伙伴约束，而没有充分恢复界面两侧的动力学分工。

本文的贡献包括三个方面：（1）构建一个同时预测 RMSF、侧链构象熵和 SASA 及其
不确定性的多链几何 Transformer，并给出从轨迹标签、局部框架到噪声感知目标的
完整定义；（2）在训练完成后引入具有三条独立 MD 重复的 ITrace 队列，以残基排序、
肽位置、接触分层和 TCR 链别动力学检验静态模型，而不是将 MD 描述本身作为模型
贡献；（3）通过同一检查点的固定坐标上下文消融，区分模型已被外部真值支持的能力
与仍待独立 apo/holo 轨迹验证的机制假设。本文不把 RMSF、$H_\chi$ 或 SASA 的变化
直接换算为结合自由能，也不把共享结合态坐标的删链差值解释为真实结合响应。
